% GENERATED by scripts/manual_tables.py from the adapter registry. Do not edit:
% rerun the script. Hand-written text: blocks/interrogation/event_families/
\ubblock{interrogation.event_families}%
  {Event families}% what the app calls it
  {Event families (SpanSet -> SpanSet + family, the Library's methods)}% registry display name
  {cluster}% insert-modal tab
  {SpanSet}% input
  {SpanSet}% output
  {Adapters/interrogation\_event\_families.py}% code
  {interrogation/event_families}% hand-written folder under blocks/
\ubblockdescription{interrogation.event_families}{EXPERIMENTAL: its template beat the shape null on 1 of the 5 real spans that held events (rec 542's known chain). On rec 6's chain its two families passed their own noise tests and were far more coherent than noise, but phase-randomised copies of the same events were placed in families as often; on 24 h M2\_aug spans it found no family at all (docs/prompts/blocks/reports/W2-evaluation.md). Read each family's noise flag and coherence before its rhythm. Sorts the events into families of similar shape with the Library's own methods (a window around each event, resampled and z-normalised, so amplitude does not enter) and writes each event's family onto it --- and, by default, as its group, so Intervals after this block gives each family its own rhythm. Each family carries the method's noise flag (better than the recording's own quiet stretches?) and its SVD entropy beside what noise scores; events in no family are -1 with their reason. Read a family's noise flag before its rhythm: an Intervals key marked \textasciitilde{}noise\_like or \textasciitilde{}untested is not a family rhythm. The noise is the span's own quiet stretches (where no event lies) whenever the chain leaves enough of them; otherwise trend-kept copies of the members, and the family says so. Tune the method's minimum family size first (hdbscan\_min\_cluster\_size). Put Event shape before it.}
\ubblockflags{interrogation.event_families}{%
  \ubflag{estimate}{}%
}
\ubblocksideinputs{interrogation.event_families}{%
}
\ubblockparams{interrogation.event_families}{%
  \ubparam{anchor}{\texttt{start}}{one of \texttt{onset}, \texttt{extremum}, \texttt{start}}{What each event's window is lined up on: its span start (default: where the detector's cut crossed, the same point for every event of one shape), Event shape's onset, or its extremum}%
  \ubparam{length\_\allowbreak{}s}{600.0}{number from 0 to 86400}{Window length in seconds. 600 (a placeholder, measured on synthetic shapes: the length that kept two shapes apart in every case of a 432-run sweep). 0 = auto, as Event windows (twice the median event duration, clamped to 60-7200 s) - fragile, warned when it clamps, overlaps or cuts the recovery}%
  \ubparam{pre\_\allowbreak{}frac}{0.2}{number from 0 to 0.95}{Share of the window before the anchor}%
  \ubparam{max\_\allowbreak{}nan\_\allowbreak{}frac}{0.1}{number from 0 to 1}{Leave out a window with more than this fraction of missing samples (fewer are interpolated)}%
  \ubparam{write\_\allowbreak{}group}{yes}{yes or no}{Write group = family, so Intervals after this block times each family on its own}%
  \ubparam{method}{\texttt{hdbscan}}{one of \texttt{hdbscan}, \texttt{kshape}, \texttt{ward}}{How families are made --- the Library's methods for shape: hdbscan (crowds of similar shapes; stragglers left out; no number of families to choose), kshape (k families compared after sliding one shape along the other), ward (a tree cut at a distance; every event placed, no noise test: untested)}%
  \ubparam{seed}{0}{whole number from 0 to 2147483647}{Seed of each family's SVD-entropy noise draws (the method's own seed is its <method>\_seed)}%
  \ubparam{hdbscan\_\allowbreak{}min\_\allowbreak{}cluster\_\allowbreak{}size}{10}{whole number from 2 to 100000}{hdbscan: minimum family size --- the smallest crowd of shapes that counts as a family --- the Library's ten-member floor. Larger finds fewer, bigger families.}%
  \ubparam{hdbscan\_\allowbreak{}min\_\allowbreak{}samples}{0}{whole number from 0 to 1000000}{hdbscan: neighbours for a core shape --- how many close neighbours a shape needs to sit inside a crowd; 0 uses the minimum family size (sklearn's default). Larger leaves more out as noise.}%
  \ubparam{hdbscan\_\allowbreak{}cluster\_\allowbreak{}selection\_\allowbreak{}method}{\texttt{eom}}{one of \texttt{eom}, \texttt{leaf}}{hdbscan: selection --- eom keeps the most persistent crowds (fewer, larger families); leaf keeps the finest ones (more, smaller, more homogeneous families).}%
  \ubparam{hdbscan\_\allowbreak{}epsilon}{0.0}{number from 0 to 1}{hdbscan: merge closer than (epsilon) --- HDBSCAN's cluster\_selection\_epsilon on the 0-1 distance scale: families closer than this are not split apart. 0 = HDBSCAN's own choice. Not the tree cut --- HDBSCAN has no tree to cut; even 0.1 merges similar shapes.}%
  \ubparam{hdbscan\_\allowbreak{}distance}{\texttt{sbd}}{one of \texttt{sbd}, \texttt{euclidean}}{hdbscan: distance --- sbd compares shapes after sliding one along the other to the best match (k-Shape's distance); euclidean compares them point by point (Ward's scale-invariant distance).}%
  \ubparam{min\_\allowbreak{}group}{10}{whole number from 1 to 100000}{every method: minimum group --- families under this many members are omitted and flagged (spec 9.10); applied by the engine after the rows not fitted are assigned.}%
  \ubparam{hdbscan\_\allowbreak{}sample}{5000}{whole number from 0 to 1000000}{hdbscan: fit on at most --- members the density fit is computed on (its distance matrix is n x n); above this many a seeded sample is fitted and the rest join the family of their nearest fitted member if they sit within its crowd.}%
  \ubparam{hdbscan\_\allowbreak{}seed}{0}{whole number from 0 to 1000000}{hdbscan: seed --- the random seed for the sample; HDBSCAN itself is deterministic.}%
  \ubparam{resample\_\allowbreak{}length}{256}{whole number from 16 to 4096}{every method: resample length --- samples every span is resampled to before z-normalising --- Ward's 256.}%
  \ubparam{hdbscan\_\allowbreak{}null\_\allowbreak{}draws}{3}{whole number from 0 to 50}{hdbscan: noise draws --- how many trend-kept surrogates of each family member (its slow trend kept, the rest phase-randomised: noise of its own colour) the family is measured against; a family is confident only if its members sit closer to its medoid than 0.9 x (a placeholder margin) the 5th percentile of those surrogates. Wherever the members' recordings leave enough quiet stretches, windows cut from them judge too (0.6 x their selection-matched quantile) and a family passes on either. 0 = no null: families are reported untested.}%
  \ubparam{kshape\_\allowbreak{}k}{0}{whole number from 0 to 1000000}{kshape: families (k) --- how many families to make. 0 chooses k from 2 to 12 by the best silhouette among the k that give two or more families of the minimum group size; every k's score is recorded with the grouping.}%
  \ubparam{kshape\_\allowbreak{}n\_\allowbreak{}init}{10}{whole number from 1 to 1000}{kshape: random starts --- k-Shape starts from a random assignment; it is run this many times and the tightest run kept. More starts, steadier families, proportionally slower.}%
  \ubparam{kshape\_\allowbreak{}seed}{0}{whole number from 0 to 1000000}{kshape: seed --- the random seed for the starts and for the sample; the same seed gives the same families.}%
  \ubparam{kshape\_\allowbreak{}omit\_\allowbreak{}d}{0}{number from 0 to 1}{kshape: further from its centre than --- 0 (the default) = automatic: a member is omitted as past\_cut when it sits further from its family's centre than 95 \% of noise does --- windows of the members' own recordings where no event lies, else trend-kept noise of the same colour. A number > 0 is a fixed shape-based distance instead (0 = a shifted copy); note that Ward's 0.5 almost never fires on this scale - two unrelated noise windows already sit near 0.42.}%
  \ubparam{kshape\_\allowbreak{}null\_\allowbreak{}draws}{3}{whole number from 0 to 50}{kshape: noise draws --- how many trend-kept surrogate sets (each window's slow trend kept, the rest phase-randomised: noise of its own colour) the families are measured against, beside the recording's own quiet stretches (which judge wherever the members' recordings leave enough of them). 0 = no null: families are reported untested and the automatic omit\_d omits nothing.}%
  \ubparam{kshape\_\allowbreak{}sample}{5000}{whole number from 0 to 1000000}{kshape: fit on at most --- members the fit is computed on; above this many a seeded sample is fitted and the rest are assigned to their nearest family centre. Recorded with the grouping.}%
  \ubparam{ward\_\allowbreak{}cut}{0.42}{number from 0 to 2}{ward: cut --- the distance to cut the tree at; the merge-height histogram shows where the tree actually merges.}%
  \ubparam{ward\_\allowbreak{}omit\_\allowbreak{}d}{0.5}{number from 0 to 1}{ward: nearest family past --- a motif whose nearest family medoid is further than this fits no family: omitted and flagged as past\_cut (spec 9.10).}%
}
